\begin{abstract}
\textbf{ABSTRACT:}
Digitalization has increasingly transformed physical infrastructures into connected cyber-physical systems (CPS). Buildings, shared facilities, rental spaces, and workplaces are no longer managed only through manual processes, but through software platforms, cloud services, Application Programming Interfaces (APIs), and connected devices. In this context, digital access control has become especially relevant, because decisions about physical entry are increasingly handled by cloud-based systems.

Simple mechanisms such as static keypad codes are easy to deploy and use, but they can be shared, leaked, or reused indefinitely without any check at the moment of entry. Stronger access control mechanisms address this by re-verifying a request at the moment of access. This leads to a security-cost trade-off: higher assurance can reduce the risk of unauthorized access, but it also increases infrastructure resource usage, operational complexity, and cloud cost.

This paper examines this trade-off using a microservice-based booking platform for music studios deployed on Amazon Web Services (AWS). The system consists of a User Service, Studio Service, Booking Service, and Access Service. Three access control scenarios with increasing security levels are implemented and compared. The first scenario generates the access code at booking time, and the code is then valid on its own with no further check at the door. The second scenario instead generates the code only at the door, after the system validates the booking code against the current time window. The third scenario builds upon the second by additionally requiring a face verification step before the code is generated, confirming that the person present at the door matches the booking.

To quantify this trade-off, the evaluation measures cost per access attempt, average Central Processing Unit (CPU) utilization, and average memory utilization. These metrics are selected because they show how additional verification logic affects the infrastructure footprint of the Access Service, which contains the scenario-specific access control logic. By comparing the three scenarios under the same deployment conditions, this work provides a practical basis for analyzing the operational cost of stronger access control in smart-lock-based CPS.
\end{abstract}